\chapter{Pecahan sebagai Bilangan}\label{ch:g5:fractions}

\omterm{def:g4:fractions:def}{Pecahan} bukan sekadar gambar sepotong pai: ia benar-benar sebuah
\emph{bilangan}, yang punya tempat pada garis bilangan --- kadang
tempat yang persis sama dengan bilangan desimal. Bab ini
menghubungkan kedua dunia itu. (Aturan lengkap hitung \omterm{def:g4:fractions:def}{pecahan}
menyusul pada \cref{ch:g6:fractions} dan seterusnya.)

\section{Pecahan pada garis, sekali lagi}

\begin{example}[Meletakkan dan membaca]\label{ex:g5:fractions:place}
Pada garis yang dipotong menjadi perlimaan, $\frac{7}{5}$ berjarak
$7$ langkah dari $0$: melewati $1$ (yang sama dengan $\frac55$),
tepatnya di $1 + \frac25$. Setiap \omterm{def:g4:fractions:def}{pecahan} tinggal di suatu tempat
pada garis; \omterm{def:g4:fractions:def}{pembilang} yang lebih besar (dengan \omterm{def:g4:fractions:def}{penyebut} yang sama)
berarti letaknya lebih ke kanan.
\end{example}

\begin{proposition}[Satu titik, banyak nama]\label{prop:g5:fractions:equal}
Memotong setiap bagian menjadi dua (atau tiga\,\dots) tidak
memindahkan titiknya:
\[
\frac{1}{2} = \frac{2}{4} = \frac{3}{6} = \frac{5}{10},
\qquad
\frac{2}{3} = \frac{4}{6} = \frac{20}{30}.
\]
Mengalikan (atau membagi) \omterm{def:g4:fractions:def}{pembilang} dan \omterm{def:g4:fractions:def}{penyebut} dengan bilangan
yang sama memberi nama lain untuk \omterm{def:g4:fractions:def}{pecahan} yang sama.
\end{proposition}
\admitted

\begin{omfigure}
\begin{tikzpicture}[scale=0.62]
  \foreach \i in {0,1} \draw[omExo] (\i*4,0) rectangle (\i*4+4,1.0);
  \fill[omDef!40] (0.06,0.06) rectangle (3.94,0.94);
  \draw[thick] (0,0) rectangle (8,1.0);
  \node[right, font=\small] at (8.3,0.5) {$\frac12$};
  \begin{scope}[yshift=-1.7cm]
  \foreach \i in {0,...,3} \draw[omExo] (\i*2,0) rectangle (\i*2+2,1.0);
  \fill[omDef!40] (0.06,0.06) rectangle (1.94,0.94);
  \fill[omDef!40] (2.06,0.06) rectangle (3.94,0.94);
  \draw[thick] (0,0) rectangle (8,1.0);
  \node[right, font=\small] at (8.3,0.5) {$\frac24$};
  \end{scope}
  \begin{scope}[yshift=-3.4cm]
  \foreach \i in {0,...,9} \draw[omExo] (\i*0.8,0) rectangle
    (\i*0.8+0.8,1.0);
  \foreach \i in {0,...,4} \fill[omDef!40] (\i*0.8+0.05,0.06)
    rectangle (\i*0.8+0.75,0.94);
  \draw[thick] (0,0) rectangle (8,1.0);
  \node[right, font=\small] at (8.3,0.5) {$\frac{5}{10}$};
  \end{scope}
\end{tikzpicture}

{\small Tiga batang, tiga nama, satu takaran: $\frac12 = \frac24 =
\frac{5}{10}$.}
\end{omfigure}

\section{Pecahan dan bilangan desimal}

\begin{proposition}[Pecahan yang punya nama desimal]\label{prop:g5:fractions:decimal}
\omterm{def:g4:fractions:def}{Pecahan} persepuluhan, perseratusan dan perseribuan \emph{adalah}
bilangan desimal:
\[
\frac{3}{10} = 0.3, \qquad \frac{27}{100} = 0.27 .
\]
Beberapa \omterm{def:g4:fractions:def}{pecahan} lain juga punya nama desimal, yang ditemukan dengan
beralih ke persepuluhan atau perseratusan:
\[
\frac{1}{2} = \frac{5}{10} = 0.5, \qquad
\frac{1}{4} = \frac{25}{100} = 0.25, \qquad
\frac{3}{4} = 0.75, \qquad
\frac{1}{5} = \frac{2}{10} = 0.2 .
\]
Tetapi tidak semuanya: $\frac13 = 0.333\dots$ tidak pernah berakhir
--- \omterm{def:g4:fractions:def}{pecahan} itulah satu-satunya nama yang tepat baginya.
\end{proposition}
\admitted

\begin{omfigure}
\begin{tikzpicture}[scale=1.55]
  \draw[-Stealth] (-0.2,0) -- (7.6,0);
  \foreach \i in {0,...,28} \draw ({\i*0.25},-0.05) -- ({\i*0.25},0.05);
  \foreach \i in {0,4,...,28} \draw ({\i*0.25},-0.09) -- ({\i*0.25},0.09);
  \foreach \x in {0,...,7} \node[below, font=\small] at (\x,-0.1) {$\x$};
  \fill[omDef] (0.5,0) circle (1.6pt)
    node[above, font=\small] {$\frac12 = 0.5$};
  \fill[omThm] (2.75,0) circle (1.6pt)
    node[above, font=\small] {$\frac{11}{4} = 2.75$};
  \fill[omMeth] (6.25,0) circle (1.6pt)
    node[above, font=\small] {$\frac{25}{4} = 6.25$};
\end{tikzpicture}

{\small Satu garis, dua bahasa: setiap titik dapat membawa nama
\omterm{def:g4:fractions:def}{pecahan} dan nama desimal.}
\end{omfigure}

\begin{example}[Memilih nama yang lebih praktis]\label{ex:g5:fractions:choose}
Untuk membayar tiga perempat dari $12$: nama \omterm{def:g4:fractions:def}{pecahan} lebih praktis,
$\frac34$ dari $12$ adalah $9$. Untuk membandingkan $\frac34$ dan
$0.8$: nama desimal lebih praktis, $0.75 < 0.8$. Menguasai dua
bahasa itu menguntungkan.
\end{example}

\section{Membandingkan pecahan}

\begin{method}[Tiga alat pembanding]\label{met:g5:fractions:compare}
\begin{enumerate}
  \item \emph{\omterm{def:g4:fractions:def}{Penyebut} sama}: yang bagiannya lebih banyak menang,
        $\frac57 > \frac37$;
  \item \emph{\omterm{def:g4:fractions:def}{pembilang} sama}: bagian yang lebih besar menang, jadi
        \omterm{def:g4:fractions:def}{penyebut} yang \emph{lebih kecil} menang: $\frac35 > \frac38$
        (perlimaan lebih besar daripada perdelapanan);
  \item \emph{patokan}: bandingkan masing-masing dengan $\frac12$
        atau dengan $1$: $\frac38 < \frac12 < \frac59$, jadi
        $\frac38 < \frac59$.
\end{enumerate}
\end{method}

\begin{example}\label{ex:g5:fractions:compare}
Siapa makan pizza lebih banyak: Ali dengan $\frac58$ atau Bela
dengan $\frac54$? \omterm{def:g4:fractions:def}{Pecahan} Bela lebih besar daripada $1$ (lima
\omterm{def:g3:division:half}{seperempat} lebih dari satu pizza utuh), \omterm{def:g4:fractions:def}{pecahan} Ali lebih kecil
daripada $1$: Bela makan lebih banyak --- bahkan lebih dari satu
pizza utuh.
\end{example}

\section{Latihan}

\begin{exercise}[$\star$]\label{exo:g5:fractions:1}
Gambarlah garis bilangan yang dipotong menjadi perlimaan dari $0$
sampai $2$, lalu letakkan: $\frac25$; \ $\frac55$; \ $\frac85$; \ $\frac{10}{5}$.
\end{exercise}

\begin{exercise}[$\star$]\label{exo:g5:fractions:2}
Lengkapilah nama yang sama nilainya:
$\frac12 = \frac{?}{8}$; \quad
$\frac34 = \frac{?}{8}$; \quad
$\frac{2}{5} = \frac{4}{?}$; \quad
$\frac{6}{9} = \frac{2}{?}$.
\end{exercise}

\begin{exercise}[$\star$]\label{exo:g5:fractions:3}
\omterm{def:g4:fractions:def}{Pecahan} mana yang menamai bilangan bulat? Sebutkan bilangan bulatnya
jika memang demikian: $\frac82$; \quad $\frac{9}{4}$; \quad
$\frac{15}{3}$; \quad $\frac{20}{5}$.
\end{exercise}

\begin{exercise}[$\star$]\label{exo:g5:fractions:4}
Tulislah sebagai bilangan desimal: $\frac{7}{10}$; \quad
$\frac{41}{100}$; \quad $\frac12$; \quad $\frac34$; \quad $\frac15$.
\end{exercise}

\begin{exercise}[$\star$]\label{exo:g5:fractions:5}
Tulislah sebagai \omterm{def:g4:fractions:def}{pecahan} (persepuluhan atau perseratusan): $0.9$;
\quad $0.13$; \quad $2.5$; \quad $0.75$ (lalu berikan nama \omterm{def:g3:division:half}{seperempatnya}).
\end{exercise}

\begin{exercise}[$\star$]\label{exo:g5:fractions:6}
Bandingkanlah, sambil menyebut alat yang dipakai
(\cref{met:g5:fractions:compare}):
\[
\frac47 \;?\; \frac67, \qquad
\frac35 \;?\; \frac37, \qquad
\frac25 \;?\; \frac{7}{12}, \qquad
\frac98 \;?\; \frac{11}{12} .
\]
\end{exercise}

\begin{exercise}[$\star$]\label{exo:g5:fractions:7}
Urutkan dengan nama desimalnya: $\frac12$; \ $0.4$; \ $\frac34$; \
$0.6$; \ $\frac15$.
\end{exercise}

\begin{exercise}[$\star$]\label{exo:g5:fractions:8}
\omterm{def:g3:division:half}{Seperempat} murid di kelas yang berisi $28$ orang bermain alat musik,
dan setengahnya berolahraga. Berapa murid masing-masing? Kelompok mana yang lebih besar?
\end{exercise}

\begin{exercise}[$\star$]\label{exo:g5:fractions:9}
Pasangkan setiap \omterm{def:g4:fractions:def}{pecahan} dengan waktunya: $\frac14$ jam, $\frac12$
jam, $\frac34$ jam, $\frac{1}{60}$ jam --- dengan $30$ menit, $45$
menit, $1$ menit, $15$ menit.
\end{exercise}

\begin{exercise}[$\star\star$]\label{exo:g5:fractions:10}
Tono meminum $\frac35$ isi botolnya yang $50$ cL; Ani meminum
$\frac12$ isi botolnya yang $60$ cL. Siapa yang minum lebih banyak?
(Hitunglah kedua takaran dalam cL --- membandingkan \omterm{def:g4:fractions:def}{pecahannya} saja
tidak cukup di sini, dan katakan mengapa.)
\end{exercise}

\begin{exercise}[$\star\star$]\label{exo:g5:fractions:11}
Berikan sebuah \omterm{def:g4:fractions:def}{pecahan} yang terletak tepat di antara $\frac12$ dan
$1$; lalu satu lagi yang tepat di antara $\frac12$ dan $\frac34$.
(Nama desimalnya mungkin membantu.)

\end{exercise}
